\begin{DoxyAuthor}{Author}
Carsten Burstedde, Lucas C. Wilcox, Tobin Isaac et.\+al. 
\end{DoxyAuthor}
\begin{DoxyCopyright}{Copyright}
GNU Lesser General Public License version 2.\+1 (or, at your option, any later version)
\end{DoxyCopyright}
The {\ttfamily sc} library provides functions and data types that can be useful in scientific computing. It is part of the {\ttfamily p4est} project, which uses it extensively. Consequently, it is also a dependency of all projects that use {\ttfamily p4est}. Some of the more important features are the following\+:


\begin{DoxyItemize}
\item The primary {\bfseries{build system}} is based on the GNU autotools. Several macro files reside in the config/ subdirectory. These are used in {\ttfamily sc} but also available to other projects for convenience. It is possible to nest packages this way, conveniently via git sumbodules.
\item CMake support has been added and is still in development.
\item The library provides {\bfseries{MPI}} wrappers named {\ttfamily sc\+\_\+mpi\+\_\+$\ast$}. With MPI enabled, they just call the corresponding MPI function. In case {\ttfamily sc} is configured without MPI (the conservative default), they implement initialize/finalize and collective calls as noops, which means that the {\ttfamily sc\+\_\+mpi\+\_\+$\ast$} calls do not have to be protected with the {\ttfamily \#ifdef SC\+\_\+\+MPI} construct in the code.
\item The library provides a logging framework that can be adapted by other packages. Multiple log levels are available, as well as options to output on just one or all MPI processes and to customize the style.
\item The library provides a set of data {\bfseries{containers}}, such as dynamically resizable arrays and memory pools, which are thin and lightweight.
\end{DoxyItemize}

To {\bfseries{build}} from a fresh checkout, call {\ttfamily ./bootstrap} first. Then proceed as described next (building from a tar distribution). The {\ttfamily bootstrap} script usually does not need to be re-\/executed!

To build the {\ttfamily sc} library from a tar distribution, use the standard procedure of the GNU autotools. The configure script is best called from a new empty directory with a relative path to the source. It takes the following options\+:


\begin{DoxyItemize}
\item {\ttfamily -\/-\/enable-\/debug} (recommended for development) lowers the log level for increased verbosity and activates the {\ttfamily SC\+\_\+\+ASSERT} macro for consistency checks.
\item {\ttfamily -\/-\/enable-\/mpi} ({\bfseries{recommended}} in general) expects a working MPI build environment, pulls in the mpi.\+h include file and activates the MPI compiler wrappers. If this option is not given, trivial wrappers for MPI routines are activated and the code is expected to be executed in serial only. Without this option, point-\/to-\/point messages to self will crash.
\item {\ttfamily -\/-\/disable-\/mpiio} may be used to avoid using {\ttfamily MPI\+\_\+\+File} based calls. The usage of {\ttfamily -\/-\/disable-\/mpiio} is deprecated and should not be used anymore. We check for all relevant file functions.
\end{DoxyItemize}

A typical development configure line looks as follows\+: \begin{quote}
{\ttfamily relative/path/to/configure CFLAGS=\char`\"{}-\/\+Wall -\/\+O0 -\/g\char`\"{} -\/-\/enable-\/mpi -\/-\/enable-\/debug} \end{quote}
A typical production {\bfseries{configure line}} looks as follows\+: \begin{quote}
{\ttfamily relative/path/to/configure CFLAGS=\char`\"{}-\/\+Wall -\/\+O2\char`\"{} -\/-\/enable-\/mpi} \end{quote}
To specify a target path for {\ttfamily make install} add a {\ttfamily -\/-\/prefix=$<$installdir$>$} option.

The {\bfseries{CMake}} build works as usual with that system. There are analogous options to enable the MPI compile and debug mode.

\begin{DoxySeeAlso}{See also}
\href{https://www.p4est.org/}{\texttt{ https\+://www.\+p4est.\+org/}} 

\href{https://www.gnu.org/licenses/licenses.html}{\texttt{ https\+://www.\+gnu.\+org/licenses/licenses.\+html}} 
\end{DoxySeeAlso}
